\section{Introduction}\label{sec:introduction}

Sequential portfolio choice can depend on how an investor arrived at the current wealth, not only on the wealth level itself.  Consider two paths from normalized wealth $1$: a rise to $1.4$ followed by a drawdown to $1.2$, and a fall to $0.8$ followed by a recovery to $1.2$.  The terminal wealth is identical, yet an adaptive reference can place the first path in the loss domain and the second in the gain domain.  Cumulative prospect theory (\CPT) provides a canonical language for such reference-dependent decisions through asymmetric valuation of gains and losses and nonlinear probability weighting \citep{kahneman1979prospect,tversky1992advances}.  In a sequential portfolio problem, these ingredients make the objective a nonlinear functional of the trajectory-outcome distribution rather than an additive one-period reward.

Recent first-order work makes this distributional objective tractable in reinforcement learning.  Foundational studies established prediction and control under \CPT \citep{prashanth2016cumulative}; most closely related, \citet{lepel2026policygradients} derive a policy-gradient theorem for finite-horizon \CPT-RL, construct an order-statistics Monte Carlo estimator, and prove convergence and sample-complexity guarantees for first-order stationarity.  Their result provides the natural fixed-reference benchmark for the present work.  The reference enters the \CPT criterion as a fixed quantity, whereas realized portfolio paths can themselves change the reference against which subsequent gains and losses are judged.

This distinction is supported by the behavioral-finance literature.  Security-trading experiments document reference-point adaptation after gains and losses, including asymmetric adjustment \citep{arkes2008reference}.  Adaptive and path-dependent references can materially change realization and portfolio decisions \citep{he2019realization,qin2022realization}, and recent multi-period portfolio analysis reaches the same conclusion under dynamic loss-aversion mechanisms \citep{gao2024multiperiod}.  Portfolio-wide gain--loss framing also changes disposition behavior \citep{an2024portfolio}, while recent evidence from Chinese retail investors shows that the information environment and investor characteristics systematically affect the disposition effect \citep{chen2025xueqiu}.  Taken together, these results motivate treating the reference as an evolving state rather than a fixed hyperparameter.

Making the reference endogenous is not merely an additional state variable.  The learner's policy changes realized wealth, realized wealth updates the reference, and the reference defines the next \CPT objective.  This feedback couples two problems that are separate under a fixed target.  Statistically, the \CPT gradient contains a nonlinear probability-weighting term evaluated at an unknown outcome distribution, so finite-budget plug-in evaluation need not be unbiased.  Dynamically, the target moves through both exogenous market variation and the learner's own state trajectory.  Convergence to one stationary point is therefore no longer the appropriate criterion; the relevant question is whether a feasible first-order policy can estimate each frozen target accurately enough to track a sequence of moving first-order conditions.

We develop \emph{Dynamic-Reference Cumulative Prospect Policy Gradient} (\DRCPTPG) around this decomposition.  For a frozen regularized \CPT target, a centered leave-one-out score statistic is combined with an independent randomized full/half correction to obtain exact unbiasedness with finite variance and finite expected trajectory cost.  For the moving target, a normalized projected update is analyzed through the projected residual and dynamic local regret (DLR), with the bound separating objective variation, replication error, and frozen-simulator error.  A factor-parametrized Dirichlet actor supplies valid long-only portfolio weights without tying the parameter dimension to the number of stocks, while asymmetric reference updating provides the behavioral state dynamics used in the portfolio instantiation.

The paper is organized around three linked questions.  First, can the nonlinear \CPT policy gradient be estimated without finite-sample bias while retaining usable variance, inference, and trajectory-budget guarantees?  Second, under what variation and simulation-error conditions does online updating keep the projected first-order residual small as the market and reference state evolve?  Third, does dynamic reference formation generate the path dependence and trading asymmetries predicted by behavioral finance, and do these distinctions remain visible in chronological A-share portfolios, held-out investor behavior, and heterogeneous recommendation interactions?  The separation is deliberate: portfolio returns are not used as evidence for estimator correctness, and behavioral fit is not used as a substitute for first-order tracking quality.

The theory answers the first two questions by establishing the regularized \CPT score-gradient identity, exact unbiasedness with finite expected cost and variance, studentized asymptotic inference, and the fixed-dimensional minimax trajectory-budget rate for the frozen-target estimator, together with projected-residual and DLR bounds for the online learner.  The experiments then test the same chain at increasing distance from the theorem objects: E1 isolates path dependence, E2 the estimator, E3 online tracking, and E4 the reference mechanism and implementation robustness.  E4 also shows that asymmetric reference adaptation is an amplifying and structuring mechanism rather than a prerequisite for online updating: the online advantage remains when adaptation is removed, while gain-dominant asymmetry strengthens the separation.  E5--E7 then move to chronological market, held-out behavioral, and calibrated interaction evidence.  In the A-share replay, \DRCPTPG attains the highest mean Sharpe ratio together with the lowest mean maximum drawdown, turnover, and transaction cost among the five policy-gradient methods; the behavioral and interaction layers provide separate evidence that the same reference representation remains informative beyond the controlled environment.

\paragraph{Contributions.}
\begin{enumerate}[leftmargin=2em]
\item \textbf{Endogenous-reference \CPT policy optimization.}  We formulate a finite-horizon portfolio objective in which the psychological reference point evolves with realized wealth.  This turns fixed-reference \CPT optimization into a path-dependent moving-target problem while retaining the fixed-reference case as a special case.
\item \textbf{Debiased gradient estimation with statistical guarantees.}  We establish the regularized \CPT score-gradient identity and construct a centered leave-one-out estimator with an independent randomized correction.  The resulting estimator is exactly unbiased, has finite expected cost and finite variance, obeys a multivariate central limit theorem, and attains the fixed-dimensional minimax rate under the same expected complete-trajectory budget.
\item \textbf{Online first-order tracking for a moving nonconvex target.}  We analyze a normalized projected update through the projected residual.  The resulting guarantee separates objective variation, replication error, and frozen-simulator error; it implies a dynamic-local-regret bound and, under a local error-bound geometry, tracking of the moving stationary set.
\item \textbf{Layered external validity without conflating estimands.}  E1--E4 test path dependence, estimator behavior, online tracking, and reference-mechanism robustness on theory-aligned objects.  A chronological 300-stock A-share replay then gives \DRCPTPG the highest mean Sharpe ratio and the lowest mean maximum drawdown, turnover, and transaction cost among the five policy-gradient methods; matched held-out behavior and calibrated interaction tests separately evaluate predictive alignment and recommendation value.
\end{enumerate}

\paragraph{Organization.}
\Cref{sec:related} positions the paper relative to \CPT optimization, adaptive-reference finance, nonstationary nonconvex learning, statistical debiasing, and recent portfolio and human-facing financial systems.  \Cref{sec:setup} defines the dynamic-reference portfolio objective and projected tracking criterion, \Cref{sec:algorithm} presents the debiased estimator and online update, and \Cref{sec:results} gives the estimator, minimax, projected-residual, DLR, and stationary-set results.  \Cref{sec:experiments} follows the E1--E7 evidence chain from mechanism and estimator tests to market, behavioral, and interaction evaluation, and \Cref{sec:discussion} synthesizes the implications and scope.  Proofs are collected in the appendix.
